\documentclass[UTF8,fontset=fandol]{ctexart}

\usepackage{amsmath,tikz}

\begin{document}

\section*{位移式的量纲检查}

图 1 中的小车从 A 向右匀加速运动。初速度为 $v_0$，加速度为 $a$，经过时间 $t$ 到达 B。有人把位移写成 $x=v_0t+at$。两项能相加吗？

\begin{center}\begin{tikzpicture}[>=stealth]
\draw (-.6,0)--(5,0);\draw[thick] (0,.12)rectangle(.8,.52);
\draw[->,blue] (.9,.46)--(2.1,.46) node[right]{$v_0$};
\draw[->] (.9,.95)--(1.7,.95) node[right]{$a$};
\node[below] at (.4,0){A};\node[below] at (4.2,0){B};
\draw[<->] (.4,-.65)--node[below]{$x$}(4.2,-.65);
\end{tikzpicture}
\\ 图 1\end{center}

图 1 中 $v_0$ 的量纲是长度除以时间，再乘时间便得到长度。

\begin{equation}[v_0t]=\frac{L}{T}\cdot T=L\label{eq:length-from-speed}\end{equation}

加速度每乘一份时间就约掉分母中的一份时间。两项量纲不同，原候选式不能给出位移。

\begin{equation}[at]=\frac{L}{T^2}\cdot T=\frac{L}{T}\label{eq:speed-from-acceleration}\end{equation}

把加速度乘时间平方，两份时间正好约掉。要写出匀加速位移，还需用运动规律确定系数 $1/2$。

\begin{equation}[at^2]=\frac{L}{T^2}\cdot T^2=L\label{eq:length-from-acceleration}\end{equation}

\end{document}
